A supervisão do processo de ensino em ambientes \textit{online} constitui um desafio central para a eficácia da aprendizagem. Nesse contexto, a análise sistemática de dados de interação emerge como uma abordagem fundamental para monitorar o engajamento discente de forma objetiva e proativa.

A Ciência de Dados Educacionais (CDE) e a Mineração de Dados Educacionais (MDE) oferecem o arcabouço teórico-metodológico para transformar registros de interação em indicadores significativos do envolvimento estudantil. Conforme \citeonline{souzaDeteccaoDesistenciaEstudantes2024}, a CDE aplica análise de dados para compreender práticas educacionais, integrando habilidades técnicas e sociais para gerar \textit{insights} sobre o processo de aprendizagem.

O monitoramento do engajamento \textit{online} permite identificar padrões comportamentais que podem sinalizar tanto a participação efetiva quanto riscos de desistência. \citeonline{tamadaPredictingStudentsRisk2022} demonstra em seu estudo que é possível extrair métricas relacionadas à assiduidade, entrega de tarefas e acesso a materiais, que refletem diretamente o envolvimento do aluno com o curso.

A detecção precoce de estudantes em risco representa uma aplicação crucial desse monitoramento. O uso de algoritmos de classificação para identificar potenciais desistências a partir de dados de interação é uma proposta de uso dessas técnicas, comparando predições com classificações reais fornecidas por docentes \cite{souzaDeteccaoDesistenciaEstudantes2024}.

Essa abordagem baseada em dados permite transitar de uma supervisão reativa para uma gestão proativa do engajamento, onde educadores podem intervir de maneira tempestiva e direcionada, potencializando a retenção e o sucesso acadêmico em ambientes de aprendizagem \textit{online}.